% =====================================================================
%  ch7_canonical.tex — §7 典范映射
% =====================================================================
\section{典范映射}\label{s:canonicalmap}

回忆：对域 $\mk$ 上亏格 $g\geq 2$ 的光滑射影曲线 $C$，有典范映射 $\operatorname{can}:C\to \PP(H^0(C,\omega_C)^*)\cong \PP(H^1(C,\cO_C))$。在 $\mk$-有理点上它可描述为
\[
p\mapsto \im\bigl(\Ext^1(\cO_p,\cO_C)\otimes\Ext^0(\cO_C,\cO_p)\to \Ext^1(\cO_C,\cO_C)=H^1(C,\cO_C)\bigr).
\]
除非 $C$ 超椭圆，此映射是闭嵌入；超椭圆时它是到其像的 $2:1$ 映射。

注意即使 $C$ 是算术亏格 $g\geq 2$ 的约化奇异曲线，仍有映射 $\operatorname{can}:C^{sm}\to \PP(H^1(C,\cO_C))$。此外，对任意 Zariski 开（相应地解析开）子集 $U\subset C^{sm}$ 与正则（相应地解析）向量场 $\theta\in H^0(U,T_U)$，有正则（相应地解析）映射 $\operatorname{can}_{\theta}:U\to H^1(C,\cO_C)$。

下面在我们的情形中对一般三价图（例如无环的情形，尽管允许有环）在 A 侧与 B 侧分别计算此映射，并看到它们一致。

\subsection{典范映射：解析框架}\label{ss:cananalytic}

为简单起见，取非阿基米德赋范域 $K$，并把 $X_K$ 看作由 $\Z[[\{q_{e}\}]]$ 到 $K$ 的标量扩张得到的曲线，其中 $q_{e}$ 送到 $\m_K$ 的某些元素。再取 Schottky 群 $\Gamma=\Gamma_{e_0/v_0}$。

在刚性解析几何框架中，$X_K$ 与商 $(\PP^1_K-F)/\Gamma$ 等同，其中 $F$ 是群 $\Gamma$ 的极限点集（等价地，$F$ 是 $\Gamma$ 的非单位元不动点集之闭包）。取 $\PP^1_K$ 上的有理函数 $\phi$，使它在 $F$ 的每一点正则。则 $\phi$ 在其极点处的主部之全体定义了一个类 $[\phi]\in H^1(X_K,\cO_K)$。

来计算这个类。注意
\[
H^1(X_K,\cO_{X_K})\cong H^1(\Gamma,K)=\Hom(\Gamma/[\Gamma,\Gamma],K).
\]
取点 $t_0\in \PP^1_K-F$，使 $\phi$ 在 $\Gamma t_0$ 的每一点正则。则有良定义的解析函数
\[
f_{\phi}(t):=\sum_{g\in\Gamma} (\phi(gt)-\phi(gt_0)),
\]
它 $\Gamma$-不变至多相差一个常数。相应的类 $[\phi]\in H^1(X_K,\cO_{X_K})=H^1(\Gamma,K)$ 由余圈
\begin{equation}
c_{\phi}(\gamma)=f_{\phi}(t)-f_{\phi}(\gamma t)=\sum_{g\in\Gamma}(\phi(g\gamma(t_0))-\phi(g(t_0))),\quad \gamma\in\Gamma
\end{equation}
给出。此余圈当然不依赖 $t_0$ 的选取。此外，若 $\gamma\neq 1$，而 $y_{\gamma}^0,y_{\gamma}^{\infty}\in \PP^1_K$ 是 $\gamma$ 的不动点（$y_{\gamma}^0$ 为吸引不动点），则
\begin{equation}
c_{\phi}(\gamma)=\sum_{\bar{g}\in\Gamma/\gamma^{\Z}}\bigl(\phi(g(y_{\gamma}^0))-\phi(g(y_{\gamma}^{\infty}))\bigr).
\end{equation}

现在若有一个解析开子集 $U\subset \PP^1_K-F$，使对所有 $g\in\Gamma\setminus\{1\}$ 有 $U\cap g(U)=\emptyset$，则 $U\cong \pr(U)\subset X_K$。在 $U$ 上取向量场 $t\frac{\partial}{\partial t}$，得到提升的典范映射 $\operatorname{can}_{t\frac{\partial}{\partial t}}:U\to H^1(\Gamma,K)$。由上述讨论，此映射由
\begin{equation}\label{eq:can_map_analytic}
\operatorname{can}_{t\frac{\partial}{\partial t}}(s)=c_{\frac{s}{t-s}}\in H^1(\Gamma,K),\qquad
c_{\frac{s}{t-s}}(\gamma)=\sum_{g\in\Gamma}\Bigl(\frac{s}{g\gamma(t_0)-s}-\frac{s}{g(t_0)-s}\Bigr)
\end{equation}
给出。

下面将在形式概形框架与 Fukaya 框架中分别看到这个 $1$-余圈如何出现。

\subsection{典范映射：形式概形}\label{ss:canformal}

这里用 $\mX$ 记前面引入的 $\Z[[\{q_{e}\}]]$ 上的形式概形，或它到某个（性质良好的）拓扑环 $R$ 的标量扩张（其中 $q_{e}$ 送到某些拓扑幂零元）。仍固定 $e_0/v_0$ 与相应的 Schottky 群 $\Gamma=\Gamma_{e_0/v_0}$。

回忆开子集 $\cU_{e},\cW_{v}\subset\mX$。注意每个交集 $\cU_{e}\cap \cU_{e'}$（$e\neq e'$）或为空，或形如 $\cW_{v}$，或形如 $\cW_{v}\sqcup\cW_{v'}$。因此对 $\mX$ 上的凝聚层 $\cF$，可把 $\cF$ 关于覆盖 $\{\cU_{e}\}$ 的 \v{C}ech 复形（拟同构地）改写成如下复形：
\[
\cK(\cF):=\Bigl\{\bigoplus_{e\in E}\Gamma(\cU_{e},\cF)\xto{d}\bigoplus_{v\in V} \Gamma(\cW_{v},\cF)\tens{\Z}V_{v}\Bigr\},
\]
其中
\[
V_{v}=(\Z\cdot e_{e_1/v}\oplus \Z\cdot e_{e_2/v}\oplus \Z\cdot e_{e_3/v})/\Z\cdot (e_{e_1/v}+e_{e_2/v}+e_{e_3/v}),
\]
而
\[
d(\{f_{e}\})_{v}=\sum_{e'/v}f_{e'}e_{e'/v}.
\]

不难直接验证有拟同构子复形 $\cK_{const}(\cO)\subset \cK(\cO)$，由（在 $\cU_{e}$ 与 $\cW_{v}$ 上的）常局部截面组成。可显式写出等同 $H^1(\cK_{const}(\cO))\cong H^1(\Gamma,R)$。即对 $e/v,e'/v$，记 $\xi^{e,e'}_{v}:V_{v}\to\Z$ 为泛函 $e_{e/v}^*-e_{e'/v}^*$。则元素 $\{a_{v}\}\in \cK^1_{const}(\cO)$ 定义了余圈
\begin{equation}\label{eq:from_Cech_to_group_cohom}
c_a\in H^1(\Gamma,R),\quad c_a(\gamma_P)=\xi_{v_1}^{e_1,e_2}(a_{v_1})+\dots+\xi_{v_n}^{e_n,e_1}(a_{v_n}),
\end{equation}
其中 $P=(v_0,e_1,v_1,\dots,e_n,v_n=v_0)$。

现在取 $\PP^1_R$ 上有理函数 $\phi$，使它在 $0,1,\infty$ 处正则。这里的意思是 $\phi(t)=\frac{h_1(t)}{h_2(t)}$，其中 $h_2(t)$ 首一、$\deg(h_1)\leq\deg(h_2)$，且 $h_2(0),h_2(1)\in R$ 可逆。于是得到凝聚层 $\cF_{h_2}\supset \cO$，使 $\supp(\cF_{h_2}/\cO)\subset \cW_{v_0}$ 且 $\cF(\cW_{v_0})=\frac{1}{h_2(t_{e_0/v_0})}\cO(\cW_{v_0})$。则 $\phi(t_{e_0/v_0})$ 的「主部」给出 $H^1(\mX,\cO_{\mX})$ 中良定义的元素。来计算此元素在 $\cK^1_{const}(\cO)$ 中的代表元。

先取截面 $f_{e}\in\Gamma(\cU_{e},\cF_{h_2})$，由
\[
f_{e}=\sum_{\substack{P=(v_0,e_1,\dots,e_n,v_n);\\ e/v_n,\,e\neq e_n}}\bigl(\phi(\gamma_P^{e_0,e}(T_{e/v_n}))-\phi(\gamma_P^{e_0,e}(0))\bigr)
\]
给出（可直接验证 $f_{e}$ 良定义）。再注意 $d(\{f_{e}\})\in \cK^1(\cF)$ 含于 $\cK^1_{const}(\cO)\subset\cK^1(\cO)\subset \cK^1(\cF)$。

于是 $d(f_{e})$ 就是所要的常值代表元，它按公式~\eqref{eq:from_Cech_to_group_cohom} 给出 $H^1(\Gamma,R)$ 中的类。直接的组合验证表明结果确与式~\eqref{eq:can_map_analytic} 一致。现在取 $s\in R$ 使 $s(1-s)$ 可逆，则类 $\operatorname{can}_{t_{e_0/v_0}\,\partial/\partial t_{e_0/v_0}}\in H^1(\mX,\cO_{\mX})$ 仍由公式~\eqref{eq:can_map_analytic} 给出。

\begin{remark}
为理解 $|s|<1$ 时的典范映射，需如注~\ref{remark:inverting_q} 那样反演 $q_{e_0}$；计算方式完全相同。
\end{remark}

\subsection{典范映射：Fukaya 范畴}\label{ss:canfukaya}

这里取如上的奇异辛流形 $M$；回忆辛面积记为 $A_{e}$，$e\in E$。仍固定 $e_0/v_0$，并取 $v_0'\neq v_0$、半边 $e_0/v_0'$。

取 v.b. 型拉格朗日 $L_0$，配以平凡秩一局部系统，它在镜像对称下对应 $\cO_{\mX}$；把其 $M_{e_0}$ 分量从 $v_0'$ 定向到 $v_0$。

Floer 复形 $\Hom(L_0,L_0)$ 就是计算图 $G$（顶点为 $v$、边为 $e$）上同调的复形。记 $p_{v}$ 与 $z_{e}$ 为其次数 $0$ 与次数 $1$ 的生成元，对应 $\cX(L_0,L_0)=L_0^+\cap L_0$ 的点。（回忆 $L_0^+=\phi^1_{\varepsilon h}(L_0)$ 是 $L_0$ 在每个顶点 $p_v$ 附近沿逆时针方向的轻微推离，它与 $L_0$ 交于各顶点，并在每个分量 $M_e$ 内还交于一点。）

现在取点型对象 $L_1$：$M_e$ 分量上的一个圆，置于 $z_{e_0}$ 与 $p_{v_0}$ 之间。取定向使 $\Hom(L_0,L_1)$ 集中在次数零，故 $\Hom(L_1,L_0)$ 集中在次数 $1$（局部系统取平凡）。令 $y_1:=L_0^+\cap L_1$，$y_2:=L_1\cap L_0$。则 $y_1\in\Hom(L_0,L_1)$、$y_2\in\Hom(L_1,L_0)$。我们关心
\[
\mu^2(y_2,y_1)=\sum_{e}a_{e}z_{e}\in\Hom^1(L_0,L_0).
\]
记 $B$ 为以 $L_1$ 为边界、含节点 $v_0$ 的那个半球的面积。

下面确定常数 $a_{e}$。先看 $e=e_0$：$L_0$、$L_1$ 与 $L_0^+$ 在 $M_{e_0}$ 内围出一个小而细的三角形，顶点为 $y_1,y_2,z_{e_0}$；相应扰动圆盘面积为零（其两条边落在 $L_0$ 上），故其面积权重为 $1$。其余全纯带都将经节点传播，为计数它们引入一些记号。

对 $e/v,e'/v$，记 $C^{v;e,e'}_{k,l}\in\Z$（$k,l\geq 0$）为满足
\[
\frac{1}{g_{v}^{e,e'}(t)^k}=\sum_{l\geq 0}C^{v;e,e'}_{k,l}t^{l}
\]
的常数。对 $k\geq 1$ 这些正是第~\ref{s:Amodel} 节引入的传播系数；常数 $C^{v;e,e'}_{0,l}=\delta_{0,l}$ 不参与传播规则，但收录进来较方便。还采用记号
\[
\delta_{v}^{e,e'}=\begin{cases}1 & \text{对 }e/v,\,e'/v,\,e\neq e';\\ 0 & \text{其他情形}\end{cases}.
\]

现在，对 $a_{e}$ 有贡献的扰动传播全纯带（除上述三角形外）分为两类：
\begin{enumerate}
\item[(I)] 先以某次数 $k>0$ 到达 $v_0$，再沿（图中的）某条路径传播，最后以次数 $0$ 到达分量 $e$ 的带；
\item[(II)] 把 $v_0$ 换成 $v_0'$ 的同类带。
\end{enumerate}

(I) 型带的贡献是如下和：
\begin{equation}\label{eq:strips_through_lambda_0}
a_{e,v_0}=\sum_{\substack{P=(v_0,e_1,v_1,\dots,e_n,v_n);\\ e/v_n,\,e\neq e_n,\,e_1\neq e_0}}\Bigl(\sum_{k,k_1,\dots,k_n>0}C^{v_0;e_0,e_1}_{k,k_1}C^{v_1;e_1,e_2}_{k_1,k_2}\dots C^{v_n;e_n,e}_{k_n,0}\,T^{\,kB+\sum_{l=1}^n k_i A_{e_i}}\Bigr).
\end{equation}
注意如下恒等式：对~\eqref{eq:strips_through_lambda_0} 中的约化路径 $P$ 与 $k>0$，
\begin{equation}\label{eq:identity_for_Schottky_groupoid}
\sum_{k_1,\dots,k_n\geq 0}C^{v_0;e_0,e_1}_{k,k_1}C^{v_1;e_1,e_2}_{k_1,k_2}\dots C^{v_n;e_n,e}_{k_n,0}q_{e_1}^{k_1}\dots q_{e_n}^{k_n}=\gamma_P^{e_0,e}(0)^{-k}
\end{equation}
（注意 $k_i$ 取非严格不等式）。若 $n>0$，记 $P'$ 为从 $P$ 去掉最后一条边所得的路径，其顶点仍从 $v_0$ 数到 $v_{n-1}$。则由~\eqref{eq:identity_for_Schottky_groupoid} 得
\begin{equation}\label{eq:identity_with_positive_k_i}
\sum_{k_1,\dots,k_n>0}C^{v_0;e_0,e_1}_{k,k_1}\dots C^{v_n;e_n,e}_{k_n,0}q_{e_1}^{k_1}\dots q_{e_n}^{k_n}
=\begin{cases}\gamma_P^{e_0,e}(0)^{-k}-\gamma_{P'}^{e_0,e_{n}}(0)^{-k} & n>0;\\
g_{v_0}^{e_0,e}(0)^{-k} & n=0.\end{cases}
\end{equation}
把~\eqref{eq:identity_with_positive_k_i} 与~\eqref{eq:strips_through_lambda_0} 结合，并把 $q_{e}$ 与 $T^{A_{e}}$ 等同，得
\begin{equation}\label{eq:strips_through_lambda_0_final}
a_{e,v_0}=\sum_{\substack{P=(v_0,e_1,v_1,\dots,e_n,v_n);\\ n>0,\,e/v_n,\,e\neq e_n,\,e_1\neq e_0}}\pm\Bigl(\frac{T^B}{\gamma_P^{e_0,e}(0)-T^B}-\frac{T^B}{\gamma_{P'}^{e_0,e_n}(0)-T^B}\Bigr)\pm\delta_{v_0}^{e_0,e}\frac{T^B}{g_{v_0}^{e_0,e}(0)-T^B}.
\end{equation}
(II) 型带完全类似。注意恒等式
\[
\frac{\frac{q}{s}}{t-\frac{q}{s}}=-\Bigl(\frac{s}{\frac{q}{t}-s}+1\Bigr),
\]
(II) 型带的贡献等于
\begin{equation}\label{eq:strips_through_lambda_0'}
a_{e,v_0'}=\sum_{\substack{P=(v_0,e_0,v_0',e_1,v_1,\dots,e_n,v_{n+1});\\ n\geq 0,\,e/v_{n+1},\,e\neq e_n}}\pm\Bigl(\frac{T^B}{\gamma_P^{e_0,e}(0)-T^B}-\frac{T^B}{\gamma_{P'}^{e_0,e_n}(0)-T^B}\Bigr).
\end{equation}
于是把~\eqref{eq:strips_through_lambda_0_final}、\eqref{eq:strips_through_lambda_0'} 与 $M_{e_0}$ 中的小三角形合起来，得
\begin{multline}
a_{e}=\sum_{\substack{P=(v_0,e_1,v_1,\dots,e_n,v_n);\\ n>0,\,e/v_n,\,e\neq e_n}}\pm\Bigl(\frac{T^B}{\gamma_P^{e_0,e}(0)-T^B}-\frac{T^B}{\gamma_{P'}^{e_0,e_n}(0)-T^B}\Bigr)\\
\pm\delta_{v_0}^{e_0,e}\frac{T^B}{g_{v_0}^{e_0,e}(0)-T^B}
\pm\delta_{e_0,e}.
\end{multline}

这就完成了 $\mu^2(y_2,y_1)\in\Hom^1(L_0,L_0)$ 的计算。为得到相应类 $c_{L_1}\in H^1(\Gamma,R)$ 在元素 $\gamma\in\Gamma$ 上的值，只需沿一条路径把 $\pm a_{e}$ 求和。与上一小节相同的组合论证表明
\[
c_{L_1}=\operatorname{can}_{t\frac{\partial}{\partial t}}(T^B),
\]
其中右端由式~\eqref{eq:can_map_analytic} 给出。于是我们看到，典范映射确实通过
\[
t_{e_0/v_0}=T^{B}\cdot(\text{单值性})
\]
把环带 $\{1>|t_{e_0/v_0}|>|T^{A_{e_0}}|\}$ 中的点与带一维局部系统的圆对应起来。
% [Tutor Note: 点拨] 这里的关键是式~\eqref{eq:identity_for_Schottky_groupoid}：传播系数之积恰是 Schottky 群胚中路径作用 $(\gamma_P^{e_0,e})^{-k}$ 的展开系数。A 侧的计数（经节点传播）与 B 侧 Čech 代表元的组合运算，实际上是同一组部分分式恒等式的两种写法。依据：原文 §7.3 与本工程 03_术语表。
